\section{Security analysis}\label{sec:security}

\subsection{Threats and mitigations}

\begin{table}[H]
\centering
\small
\caption{Threats and protocol-level mitigations.}
\label{tab:threats}
\begin{tabularx}{\linewidth}{@{}p{0.29\linewidth}X@{}}
\toprule
\textbf{Threat} & \textbf{Mitigation} \\
\midrule
Forge a crossing & The chosen cross leaf requires an owner signature and a
canonical 55-byte body. The body binds the sat, mint txid, mint index and
direction; the state output is a one-time UTXO. \\
Replay a witness on another frog & The body and the parent tag both name the
mint; the state output commits to the owner. A copied witness fails
\code{CHECKSIG}, the body check, or the spent-leaf check. \\
Replay a witness with a different direction & The direction is the last byte
of the body pushed by the leaf, and the leaf hash is committed by the control
block. A witness that spent the left leaf cannot authorize a straight or
right body. \\
Counterfeit a mint & A mint requires exactly one inscription reveal in input
0, no pointer, the correct content type, the canonical v8 body digest, and a
first sat that is tracked and frog-free. The v7 digest is accepted only for
the grandfathered sat; v6 and every other v7 mint are rejected. \\
Hide a death & The fold derives the wrap from $H(T+1)$ and the direction in
the body. The static band never prints a count it cannot know. \\
Choose the outcome & The author chooses $T$ and the direction, not
$H(T+1)$. The target is frozen at $T$; the wrap does not exist when the
crossing is signed. A block producer can influence the candidate hash or
withhold a block, but that is miner influence, not player agency. \\
Blame another Hopper for a block's wrap & The wrap is global and depends only
on the block hash. One frog's death says nothing about another frog's sat or
owner. \\
Kill without publishing & A transfer publishes nothing, so a frog cannot be
killed without publishing a crossing. A burn is a separate, observable
terminal event. \\
Pointer attack & HOPPER never writes a pointer tag, and a mint carrying one
is rejected. The inscription stays on the anchor sat. \\
External-origin phishing & The mint document contains no external URL. A
viewer with different recursion semantics sees the static band, not a
third-party page. \\
Accidental burn & The wallet never uses the frog address for funding and
refuses to spend a coin whose ranges it cannot attribute. \\
Parameter confusion & The parameter hash is committed in snapshots and
checked by the wallet before mint or hop. A mismatch is a hard error. A
\code{generation_len} that does not match its formula is rejected. \\
Indexer lies & The indexer is not in the art's trust path. The viewer's own
ord node supplies the satpoints, bodies, raw transactions and block hashes
used by the fold. \\
Fabricated child & A body is not believed because it names the right parent.
The fold checks the child's sat and satpoint, then proves from the raw
transaction that input 0 spent the canonical state leaf for the body's
direction. A body pasted into a funding input, or an ordinary output-only
deposit, fails the proof and renders the fold unverified. \\
Reorg & Timers are pure functions of height and hash. The indexer replays to
the fork point; the fold simply asks again. \\
\bottomrule
\end{tabularx}
\end{table}

\subsection{Properties}

The following statements are informal but precise enough to test.

\begin{enumerate}[leftmargin=1.7em,itemsep=0.35em]
  \item \textbf{Unforgeability.} Under the security of BIP-340 signatures on
        secp256k1, an adversary without the owner key cannot produce a
        witness that passes a cross or send leaf with a valid sighash.
  \item \textbf{Direction binding.} The direction is part of the body, the
        body is pushed by the leaf, and the leaf hash is committed by the
        control block and the output key. A crossing cannot be re-labelled
        with another direction, and a witness cannot be replayed with one.
  \item \textbf{No replay across frogs.} A state body names
        $(\mathrm{sat},\mathit{mint\_txid},\mathit{mint\_index})$, and the
        state output is spent once. A witness for one frog cannot satisfy
        another frog's state output or body check.
  \item \textbf{Proof-checked public verifiability.} The fold reconstructs
        state from the mint document, the crossing inscriptions, the raw
        transactions, the tip and block hashes. A body-valid child without a
        valid canonical spend makes the fold unverified rather than
        optimistic. No secret is required and no HOPPER-specific server is
        authoritative.
  \item \textbf{Outcome independence.} The verdict depends on $H(T+1)$, which
        is not known when the crossing is signed, and on the direction the
        owner committed. The owner can time the crossing and pick a
        direction; it cannot select a favorable hash.
  \item \textbf{Deterministic replay.} The state machine is a pure function
        of blocks, seeds and parameters. Two honest indexers with the same
        inputs produce the same state root and event log.
  \item \textbf{Fail-closed art.} If the fold cannot fetch a child, verify a
        satpoint or raw proof, complete pagination, or run JavaScript, the
        document presents the static band or the unverified band with an
        unknown count. It does not claim the frog is alive, and it does not
        claim \code{HOME}.
\end{enumerate}

The first three properties are the custody argument. The fourth is
proof-checking plus the game argument. The fifth and sixth are the
infrastructure arguments. The seventh is the presentation argument.

\subsection{Limitations and honest assumptions}

\begin{warnbox}{What HOPPER does not promise}
\begin{itemize}[leftmargin=1.2em,itemsep=0.25em]
  \item \textbf{Liveness.} A transaction can be censored or fee-starved. The
        protocol decides what happens after confirmation, not when
        confirmation occurs.
  \item \textbf{Privacy.} Addresses and sat histories are public. The design
        is self-custodial, not anonymous. The global wrap also correlates
        every Hopper in a block: one block-wide coin flip, not one per frog.
  \item \textbf{Data availability.} The art, state bodies and raw
        transactions must remain retrievable from the ord node or another
        index. A viewer that cannot fetch them shows the fallback or
        unverified band.
  \item \textbf{Legacy reveals.} A keyless commit made before the signed
        commit rule remains revealable forever. No new commit may use that
        form.
  \item \textbf{Discovery.} A wallet that cannot see the chain or a trusted
        index may not find a frog sent to it. The frog still exists; the
        wallet simply has not discovered it.
  \item \textbf{Rule changes.} The rules are fixed for wire version
        \code{0x08}. A change requires a new wire version; it cannot be
        deployed by editing an indexer's configuration. The single
        grandfathered v7 Hopper is a custody exception, not a general
        compatibility mode.
  \item \textbf{Wrap bias.} The modulo is biased toward left by one part in
        $2^{32}$. This is negligible for the game and is part of the
        consensus definition.
\end{itemize}
\end{warnbox}

\subsection{Why the extra transaction is unavoidable}

A state leaf names the mint inscription id. The id contains the mint reveal
txid. An output created by the reveal transaction cannot contain that txid,
because a transaction cannot commit to its own hash. The first state output
must therefore be created by a later transaction: either the first transfer
or the one-time \textsc{prepare}. This is not an implementation artifact; it
is a consequence of self-referential identity on a hash-linked ledger.

The alternative --- omitting the mint id from the state leaf --- would make
the custody script independent of the frog's identity and would weaken
replay resistance. HOPPER chooses the extra transaction and keeps the binding.

\subsection{Compatibility notes}

Three compatibility rules are normative and easy to misread.

\begin{enumerate}[leftmargin=1.6em,itemsep=0.25em]
  \item Version 8 is the current deployment. The envelope version
        (\code{0x08}), the 55-byte body, the four-leaf tree, the proof-checked
        fold and the art all changed. A v8 indexer accepts exactly one v8
        mint-body digest. For \code{GRANDFATHERED_SAT} only, the released v7
        envelope, digest, \code{V7} body and \code{FRG7CROSS}/\code{0x07}
        state path remain valid. Every other v7 mint or crossing is rejected,
        and v6 bodies are rejected. The v7 card's embedded fold is immutable
        and predates the proof check; it is not an authoritative art oracle.
  \item The indexer accepts both minimal and PUSHDATA1 spellings of short
        pushes in an ord envelope. The reference encoder emits minimal
        pushes, and byte-exact fixtures use them, but the parser is lenient
        in the same places ord is lenient.
  \item The wrap must be computed with exactly
        \code{u32LE(H(B)_internal[0..4]) mod 3}. A wider word, a rejection
        scheme, or a hash of the hash changes the outcome and is not v8.
\end{enumerate}

These rules do not change the game. They keep old bytes valid where they
must be, and they make the decoder match the ecosystem it lives in.
